% TOP_PROBLEM: {"id":30007611,"problem_number":"LOCAL-30007611","title":"CBDCs and bank funding","category_id":21,"category":{"id":21,"name":"economics","display_name":"Economics","description":"Open economic theory statements, published research agendas, and proposed scoped specifications; question provenance and historical priority are qualified per record.","slug":"economics","order_index":21,"created_at":"2026-10-08T00:00:00Z"},"status":"research_question","external_url":null,"legacy_ids":[],"published":false,"statement":"In the explicitly specified setting: How would retail CBDC adoption change ordinary bank funding, credit provision and crisis-time deposit flight under alternative designs? The mathematical statement above fixes the target, assumptions and scope of this catalogue formulation.","economics":{"original_id":100,"field":"Banking, credit and financial stability","kind":"design","formalization_basis":"adapted_research_question","source_status":"research_agenda","priority_status":"earliest_found","priority_note":"This is the earliest explicit statement verified in this audit; first priority for the full catalogue formulation is not established.","earliest_verified_reference":{"authors":["Bank of England"],"title":"Bank of England Agenda for Research, 2025–2028","year":2026,"url":"https://www.bankofengland.co.uk/research/bank-of-england-agenda-for-research","doi":null,"locator":"Financial System / Priority topics for 2026, item 3; Central banking toolkit, paragraph 2","evidence_summary":"Requests digital-asset effects and CBDC design analysis."},"current_reference":{"authors":["Bank of England"],"title":"Bank of England Agenda for Research, 2025–2028","year":2026,"url":"https://www.bankofengland.co.uk/research/bank-of-england-agenda-for-research","doi":null,"locator":"Financial System / Priority topics for 2026, item 3; Central banking toolkit, paragraph 2","evidence_summary":"Requests digital-asset effects and CBDC design analysis."},"formalization_note":"The agenda supports a broad CBDC-design research question. Some bank-funding and holding-cap models are already solved; this proposed target concerns adoption and external validity."},"statusQualification":"Published question or research agenda verified at the cited locator. The mathematical specification is adapted for this catalogue and includes additional assumptions; source publication does not establish first-ever priority or certify this exact model remains unsolved. The agenda supports a broad CBDC-design research question. Some bank-funding and holding-cap models are already solved; this proposed target concerns adoption and external validity.","statusReviewedAt":"2026-10-08"}
% FIRST_POSTED: 2026-10-08
% ENTRY_KIND: statement_only
% !TeX program = lualatex
\documentclass[11pt]{article}
\usepackage{fontspec}
\usepackage[a4paper,margin=25mm]{geometry}
\usepackage{amsmath,amssymb,mathtools}
\usepackage{xurl}
\usepackage[hidelinks]{hyperref}
\newcommand{\sref}[2]{\hyperlink{source:#1:#2}{[S#2]}}
\setlength{\emergencystretch}{3em}
\begin{document}
% problemId:30007611
\section*{CBDCs and bank funding}\label{30007611}
\subsection{Definitions and mathematical statement}
Fix households choosing among bank deposits, cash and a retail central-bank digital currency. Banks make loans financed by deposits, wholesale funding and equity. A digital-currency design \(a=(\bar q,r,\kappa)\) specifies holding cap \(\bar q\), remuneration rule \(r\) and conversion friction \(\kappa\). Households respond to payment convenience, return, privacy and perceived default risk; banks adjust deposit pricing and lending. Let \(D(a,x)\), \(L(a,x)\) and \(R(a,x)\) be equilibrium deposits, productive credit and a declared crisis-run statistic in state \(x\). The task is to identify or bound the design-dependent changes
\[(D(a,x)-D(0,x),\ L(a,x)-L(0,x),\ R(a,x)-R(0,x)).\]
The baseline \(0\) is the otherwise identical economy without retail digital currency. Distinguish slow portfolio substitution in ordinary states from rapid crisis conversion. Specify whether the central bank replaces lost bank funding and on what collateral terms, since this changes credit effects. Estimate adoption and conversion behavior using justified preference and information assumptions, and characterize welfare-maximizing caps or rates over a feasible set. A result under one bank-funding replacement assumption or one calibrated holding cap does not establish a general disintermediation effect.

\par\textbf{Formulation and scope.} The agenda supports a broad CBDC-design research question. Some bank-funding and holding-cap models are already solved; this proposed target concerns adoption and external validity.

\par\textbf{Open-question provenance.} An explicit question or research agenda was verified at the source locator. This does not by itself certify that every later version remains unresolved \sref{30007611}{1}.

\par\textbf{Historical priority.} This is the earliest explicit statement verified in this audit; first priority for the full catalogue formulation is not established.

\subsection{Short English statement}
In the explicitly specified setting: How would retail CBDC adoption change ordinary bank funding, credit provision and crisis-time deposit flight under alternative designs? The mathematical statement above fixes the target, assumptions and scope of this catalogue formulation.

\subsection{Sources}
\begin{itemize}\raggedright
\item[\textnormal{[S1]}]\hypertarget{source:30007611:1}{} Bank of England. 2026. Bank of England Agenda for Research, 2025–2028. Financial System / Priority topics for 2026, item 3; Central banking toolkit, paragraph 2.
\url{https://www.bankofengland.co.uk/research/bank-of-england-agenda-for-research}.
{\small\textit{Earliest verified question/agenda reference; historical priority qualified below. Requests digital-asset effects and CBDC design analysis.}}
\end{itemize}
\end{document}
